% Task 2 report; build from this folder with: tectonic task2_report.tex
\documentclass[11pt,a4paper]{article}
\usepackage[margin=2cm,top=1.4cm]{geometry}
\usepackage{enumitem}
\usepackage{hyperref}

% title block: left-aligned title, then author on the left and date on the right in one line
\makeatletter
\renewcommand{\maketitle}{%
  \noindent{\LARGE \@title\par}\vspace{0.6em}
  \noindent{\large \@author\hfill\@date\par}\vspace{1.5em}}
\makeatother

\title{Task 2: Reproducing the Baselines}
\author{L335 MiniProject}
\date{\today}

\begin{document}
\maketitle

% task description, copied verbatim from the project brief
% dashes at 1em, item text at 2.5em; no extra space before, between or after the items
\noindent\textbf{Task description:} 
\begin{itemize}[label=--, leftmargin=2.5em, rightmargin=2.5em, labelindent=1em, labelwidth=1em, labelsep=0.5em, align=left, topsep=0pt, partopsep=0pt, itemsep=0pt, parsep=0pt]
\item Run PGD and the geometry-only object model of P2P-Bridge using their released checkpoints on PU-Net evaluation data.
\item First reproduce each method's documented settings.
\item Then compare them using the same test shapes, point counts, noise realisations, normalisation, and metric definitions.
\item Analyse their failure cases before modifying either model.
\end{itemize}
\noindent\hspace*{1em}\textit{Note:} Full retraining of the baselines is not required.
\medskip



\end{document}
